% ============================================================================
% Chapter 6: Outlook (Headers Only)
% ============================================================================

\newpage
\section{Outlook}\label{sec:outlook}

\subsection{Conclusion}\label{sec:conclusion}

\textit{[To be written. Summary of what was developed: (1) speedrunning benchmark, (2) custom VLM-agent architecture with powerful abstractions for embodied settings, (3) AutoEvolve algorithm for self-adaptation of model harnesses via experiential learning. What was shown and why it is promising.]}

\subsection{Future Work}\label{sec:future_work}

\tersoo{Forward-reference notes from Chapter 4:
 - LLM+RL hybrid scalability: SPD-style approaches where LLMs provide task structure priors
   and RL refines execution; how does this scale beyond the tutorial to long-horizon settings?
 - SLAM/navigation without hand-crafted pathfinding: learned spatial representations,
   SLAM-like map construction from partial observations, reducing dependence on porymap data.
 - Domain-tailored abstractions for long-horizon RPG play: general-purpose coding harnesses
   (Claude Code, Gemini CLI) show promise but lack the operational granularity needed for
   reliable navigation and multi-phase quest management.}

\subsubsection{Harness Isolation and Ablation}\label{sec:harness_isolation}

\textit{[To be written. As noted in Section~\ref{sec:harness_spectrum}, $\mathcal{H}_{\min}$ and $\mathcal{H}_{\text{expert}}$ share primitive access to subagents, memory, and skills; they differ in what is pre-built rather than in fundamental capability. Developing a more strictly isolated minimal baseline, where runtime accumulation of these components is disabled, would enable cleaner ablation studies and sharpen the attribution of performance to human-engineered initialization versus automated discovery.]}

\subsubsection{Model Distillation}\label{sec:distillation}

\textit{[To be written. Training open-source models on a harness learned by a more powerful model; teacher-model paradigm via AutoEvolve; implications for industry.]}

\subsubsection{Full Game Evaluations}\label{sec:full_game}

\textit{[To be written. Extending evaluations beyond three gyms to full game completion.]}

\subsubsection{Extensions to Other Domains}\label{sec:extensions}

\textit{[To be written. Generalization of AutoEvolve to coding, robotics, and other embodied/agentic domains.]}

\subsubsection{Meta-Harness Research}\label{sec:meta_harness_future}

\textit{[To be written. The emerging field of automated harness construction represents a parallel line of research that complements and extends the contributions of this thesis. The Darwin G\"{o}del Machine (DGM)~\citep{zhang2025dgm} demonstrated open-ended evolution of self-improving coding agents through quality-diversity optimization, improving SWE-bench performance from 20.0\% to 50.0\%. Meta-Harness~\citep{lee2026metaharness} automated harness optimization by providing a coding agent with up to 10 million tokens of diagnostic context per optimization step, demonstrating a $6\times$ performance gap through harness optimization alone. These approaches operate across independent evaluation episodes and in textual environments; extending them to the online, embodied, reset-free regime explored by AutoEvolve is a natural direction for future work.]}

\subsubsection{Recursive Language Model Directions}\label{sec:rlm_future}

\textit{[To be written. The Recursive Language Model (RLM) framework~\citep{zhang2025rlm} provides a theoretical foundation for agents that programmatically interact with and restructure their own context during inference. The structured trajectory window in our expert harness and the segment-level optimization in AutoEvolve can both be understood as instances of recursive context management. Formalizing these connections and developing RLM-aware training objectives that explicitly optimize for context-management quality is a promising direction that could yield agents with stronger long-horizon coherence.]}
